\section{TUJUAN DAN MANFAAT}

\subsection{Tujuan}

Savora dirancang dengan tujuan yang spesifik, terukur, dapat dicapai, relevan, dan berbatas waktu (SMART) sebagai berikut:

\begin{enumerate}[leftmargin=*]
    \item \textbf{Specific (Spesifik)} : Mengembangkan platform manajemen \textit{food waste} berbasis web yang menghubungkan UMKM kuliner dengan konsumen dan mitra donasi untuk menyalurkan makanan surplus yang masih layak konsumsi.
    
    \item \textbf{Measurable (Terukur)} : Mengurangi volume \textit{food waste} pada UMKM kuliner mitra sebesar 40\% dalam 6 bulan pertama operasional, serta meningkatkan pendapatan UMKM dari penjualan makanan surplus sebesar 25\%.
    
    \item \textbf{Achievable (Dapat Dicapai)} : Menerapkan fitur \textit{Food Trust Index}, \textit{Food Score decay}, \textit{keyword classification}, dan pembayaran \textit{cashless} yang selaras dengan studi literatur \cite{kusumawati2025mobile}\cite{qlearning2026pricing} dan best practices platform serupa.
    
    \item \textbf{Relevant (Relevan)} : Menjawab kebutuhan mendesak akan solusi \textit{food waste} yang terintegrasi \cite{mobilefoodwaste2024review}, sesuai dengan target pengurangan sampah makanan nasional sebesar 50\% pada tahun 2030 \cite{unsdg2024}.
    
    \item \textbf{Time-bound (Berbatas Waktu)} : Mencapai target implementasi di 100 UMKM kuliner di Kota Bandung dalam 12 bulan pertama, dengan ekspansi ke kota lain pada tahun kedua.
\end{enumerate}

\subsection{Manfaat}

Manfaat Savora dapat dilihat dari tiga perspektif utama:

\subsubsection{Manfaat bagi UMKM Kuliner}
\begin{itemize}[leftmargin=*]
    \item Mengurangi kerugian akibat makanan surplus yang tidak terjual
    \item Mendapatkan pendapatan tambahan melalui penjualan makanan surplus dengan harga diskon
    \item Meningkatkan efisiensi operasional melalui prediksi produksi yang akurat
    \item Meningkatkan reputasi sebagai UMKM yang peduli lingkungan
    \item Mendapatkan akses ke jaringan donasi untuk makanan yang tidak terjual
\end{itemize}

\subsubsection{Manfaat bagi Konsumen}
\begin{itemize}[leftmargin=*]
    \item Memperoleh akses terhadap makanan berkualitas dengan harga yang lebih terjangkau
    \item Berkontribusi langsung dalam pengurangan \textit{food waste} dan emisi gas rumah kaca
    \item Mendapatkan informasi transparan mengenai kelayakan dan asal-usul makanan
    \item Memperoleh penghargaan berupa reward dan badge atas kontribusi lingkungan
\end{itemize}

\subsubsection{Manfaat bagi Lingkungan}
\begin{itemize}[leftmargin=*]
    \item Mengurangi volume sampah organik yang berakhir di TPA
    \item Menurunkan emisi gas rumah kaca dari pembusukan sampah makanan
    \item Melestarikan sumber daya alam yang digunakan dalam produksi makanan
    \item Mendorong terciptanya budaya pengurangan \textit{food waste} di masyarakat
\end{itemize}

\subsection{Dampak Jangka Panjang}

\subsubsection{Dampak Ekonomi}
Dalam jangka panjang, Savora berpotensi menciptakan ekosistem ekonomi sirkular di sektor kuliner, di mana setiap produk makanan memiliki nilai ekonomis hingga akhir siklus hidupnya. Hal ini tidak hanya menguntungkan UMKM dan konsumen, tetapi juga membuka peluang kerja baru di bidang logistik dan manajemen \textit{food waste}.

\subsubsection{Dampak Sosial}
Platform ini berpotensi memperkuat hubungan antara UMKM dengan komunitas sekitar melalui mekanisme donasi. Makanan surplus yang didonasikan dapat membantu memenuhi kebutuhan pangan kelompok masyarakat yang membutuhkan, sehingga berkontribusi pada penurunan angka ketahanan pangan.

\subsubsection{Dampak Lingkungan}
Dengan target pengurangan \textit{food waste} sebesar 40\%, Savora berkontribusi pada pengurangan emisi gas rumah kaca dan pelestarian sumber daya alam. Skalabilitas platform ini memungkinkan dampak lingkungan yang lebih luas saat diterapkan di berbagai kota di Indonesia.

\subsection{Metrik Dampak}

\begin{table}[htbp]
\centering
\caption{Metrik Dampak Savora}
\label{tab:metrik_dampak}
\begin{tabularx}{\textwidth}{|l|X|X|}
\hline
\textbf{Metrik} & \textbf{Sebelum Penerapan} & \textbf{Sesudah Penerapan} \\
\hline
Jumlah UMKM mitra & 0 UMKM & Target 100 UMKM \\
\hline
Volume food waste terkurangi & Baseline 0\% & Target 40\% \\
\hline
Pendapatan tambahan UMKM & Baseline 0\% & Target 25\% \\
\hline
Jumlah makanan diselamatkan & 0 kg & Target 15.000 kg \\
\hline
Pengurangan emisi CO$_2$ & 0 ton & Target 30 ton \\
\hline
Jumlah konsumen aktif & 0 & Target 5.000 \\
\hline
Jumlah donasi tersalurkan & 0 kg & Target 5.000 kg \\
\hline
\end{tabularx}
\end{table}
